\section{v0 y AI-native development: las restricciones de producto generan problemas de investigación}

v0 convierte ideas en lenguaje natural en un frontend ejecutable y se ha ido ampliando hasta abarcar full-stack application. La versión temprana vigente en el momento de la clase usaba GPT-3.5 Turbo; el function calling nativo todavía no estaba maduro, así que el equipo tuvo que construir una bespoke tool invocation. Esta historia muestra que un producto de IA no consiste simplemente en llamar a la API de un modelo, sino que es un sistema conjunto de modelo, herramientas, framework knowledge, preview deployment y retroalimentación de los usuarios.

\subsection{Del prompt a la aplicación en ejecución}

La discusión anterior sobre compute-density se centraba en cómo ejecutar las solicitudes de forma eficiente; aquí el enfoque pasa al ciclo cerrado de generación y verificación de una AI application. Esta sección no responde a «si el modelo puede escribir un fragmento de JSX», sino a cómo un prompt, mediante tool invocation, framework-aware code generation, build, preview y human feedback, se convierte en un producto interactivo. Al leer la figura, conviene ver el preview environment como un verifier y no solo como el punto final de la publicación.

\lecturefigure{14-v0-loop.png}{v0 conecta el código generado, las llamadas a herramientas, la preview URL y la retroalimentación de los usuarios en un ciclo cerrado de producto.}{Subtítulos locales 30:15--31:55; rediseño conceptual.}

\begin{knowledgebox}{Lectura de la figura: por qué el entorno de ejecución es el verifier del coding agent}
Que el código parezca razonable no significa que pueda hacer build, renderizarse y permitir interacción. El preview deployment ofrece retroalimentación ejecutable: si la compilación se completa, si la página es accesible y si el aspecto visual y el comportamiento cumplen los requisitos. El agent necesita devolver el build log, el browser result y las modificaciones del usuario al ciclo de generación.
\end{knowledgebox}

\begin{warningbox}{Fast growth no es architecture proof}
En la clase se mencionó el rápido crecimiento de v0, pero eso solo indica una señal de demanda fuerte; no demuestra que la economía unitaria, la retention a largo plazo o la calidad de generación estén resueltas. Un producto de IA sigue necesitando métricas de largo plazo como la distribución de tareas, la tasa de éxito, el número de rondas de corrección, el token cost, la latency y la retención de usuarios.
\end{warningbox}

\teachervoice{Rauch trata «fijar estándares exageradamente altos» como un método de investigación: si el objetivo es solo automatizar un workflow existente, es posible que el equipo nunca se tope con un novel problem; si se exige que cada prompt se convierta rápidamente en un sistema ejecutable, el tool calling, el sandbox, el preview y la retroalimentación se convierten en una nueva infraestructura de investigación.}

\subsection{Resumen de la sección}

La unidad del sistema de v0 no es una completion, sino el ciclo prompt--tool--code--preview--feedback. Las restricciones de producto pueden crear problemas de investigación, pero deben validarse con evidencia de ejecución y no solo mostrando texto generado.
